\section{Conclusion \& Future Work} \label{sec:conclusion-and-future-work}

In this Bachelor's thesis, the potential of \acp{llm} for educational question generation was investigated through two experiments: 

\begin{enumerate}
  \item Evaluating the adherence of \ac{llm}-generated questions to diverse instructional texts, and
  \item Analyzing how three prompting variants (specifying only the question format, only a targeted Bloom level, or both together) affect the generation of questions across Bloom's revised Taxonomy and their adherence to the intended cognitive level, if specified.
\end{enumerate}

Four popular \acp{llm} with \textit{reasoning} capabilities were employed to conduct \ac{aqg}: Anthropic Claude 3.7 Sonnet, DeepSeek R1, Google Gemini 2.5 Flash, and OpenAI o3.

\subsection{Key Findings}

Firstly, the quantitative results demonstrated that, based on the \textit{Adherence Score} -- an alternative to semantic similarity --, \acp{llm} can generate questions with high content fidelity across diverse instructional sources. However, the cosine similarity, between sentence transformer based embedding vectors, failed to distinguish between source-based and no-source-based questions, with similar results for both types. Sentence transformers are not suitable for the texts used in this thesis, as these models are not designed for longer educational texts, but rather for sentence-level comparisons. 

Significantly, the prompt complexity played a crucial role in both quantitative and qualitative analysis at question quality outcomes, where the common prompt consistently outperformed the complex prompt in most quality dimensions, surprisingly. The complex prompt only increased the \textit{Challenging} ratings, but at a noticeable quality cost, which was not desired. This implies that a vastly formulated prompt does not necessarily lead to better results.

Moreover, the source materials used in Experiment 1a had an impact on the question quality. The lengthy transcript was rated highest in multiple criteria, such as \textit{Clarity}, \textit{Answerability}, and \textit{Correctness}, but the lowest in \textit{Challenging}. The concise script also yielded high scoring, especially at the staff's evaluation, and achieving the highest \textit{Total Score} due to a better balance in \textit{Challenging}, while the Tanenbaum excerpts were the weakest overall, especially in \textit{Language} and \textit{Answerability}.

Model differentiation was also observed, with OpenAI and DeepSeek demonstrating higher consistency, while Anthropic and Google were more variable, particularly showing a tendency to produce linguistically complex questions and occasionally incorrect or invented answer possibilities. The \textit{Challenging} criterion was also rated lower for all models, with DeepSeek and OpenAI being even lower than Anthropic and Google. 

Regarding the \textit{Manipulation Handling} in Experiment 1b, the quantitative analysis showed that OpenAI and DeepSeek had lower \textit{Adherence Scores} when the complex prompt was used, so these models could partially detect the manipulations in the source material. In contrast, Anthropic and Google often reproduced the manipulations without any remark. The small sample size of 16/56 questions in Experiment 1b and the low inter-rater reliability shows inconsistencies in the results.

The second experiment showed that specifying Bloom levels in the prompts resulted in high adherence to the intended cognitive levels, which is crucial for effective educational assessment. Without Bloom specification, the questions tended to be at lower cognitive levels (at \textit{Understanding}\,/\,\textit{Applying}), which is not sufficient for demanding educational content. Open-Ended questions can be used across all Bloom levels and achieved superior Bloom alignment, while \acp{mcq} were more suitable for lower cognitive levels. 

Moreover, the models overall performed similarly, while being rated with almost the same experimental criteria. When considering the \textit{Bloom's Level} scores, without Bloom specification, Anthropic and Google were rated higher than DeepSeek and OpenAI, but when Bloom levels were specified, OpenAI performed best in the Bloom Only setting, while Google performed best in the Type+Bloom setting. OpenAI was occasionally underperforming in the Type+Bloom setup, which was due to the \acp{mcq} being rated lower, since Bloom level 6 contradicts the \ac{mcq} format, which was not considered in the prompt design.

Across the experiments, linguistic complexity issues were observed, where overly complex questions and often occurring phrases were reducing the \textit{Language} and \textit{Value} ratings. Furthermore, the rubric designs of the experiments reduced the agreements due to its broad 0-10 scale, coupled with not sufficiently clear anchor descriptions. Unfortunately, the \textit{Value} criterion was criticized, since in this thesis, it was not planned to integrate explicit learning objectives for the questions, which made the evaluation more abstract and less clear.

All in all, the \textit{reasoning}-based \acp{llm} can automate the generation of relevant, clear, and Bloom-aligned questions, but show diversity in criteria such as \textit{Correctness}, \textit{Challenging}, and especially in \textit{Clarity}\,/\,\textit{Language}. Based on the findings, questions with high content fidelity, and proper linguistic quality across the Bloom levels require: (1) Concise, well-structured source inputs rather than overly complex ones; (2) Explicit revised Bloom's Taxonomy level explanations -- best coupled with learning objectives for each question\,/\,topic; (3) Proper but restrained prompt complexity with clear phrasing; and (4) Refined evaluation rubrics.

\subsection{Answering the Research Questions}

\requ{1}{To what extent do \acp{llm} adhere to the content of diverse provided instructional texts when generating questions?}

The quantitative analysis yielded promising results, with high \textit{Adherence Score} entries across all models and prompts. However, cosine similarity was unsuitable for this task, as it failed to distinguish between source-based and no-source-based questions, since sentence transformers fail with longer educational texts. The common prompt consistently outperformed the complex prompt in content adherence across both quantitative measures and expert evaluations. Taking the qualitative analysis into account, the mentioned performance gap between the common and complex prompt was also visible in most criteria, where structured materials (script and transcript) were mostly rated higher than the Tanenbaum book excerpts.

Moreover, the models showed different strengths and weaknesses, with OpenAI and DeepSeek being more consistent in \textit{Correctness}, generating questions that well reflect the source material without unnecessary linguistic complexity. Anthropic and Google were more variable, often generating questions which were linguistically complex (especially in the complex prompt) and occasionally incorrect or invented answer possibilities. Also noteworthy, the DeepSeek and Google models both struggled most with proper question phrasing, since these generated many questions including \enquote{based on the text} formulations. 

The quantitative analysis also showed that the models can detect manipulations in the source material, shown in a lower \textit{Adherence Score} for OpenAI and especially DeepSeek -- coupled with the complex prompt --, while Anthropic and Google more often reproduced the manipulations without any remark. Qualitatively, the \textit{Manipulation Handling} criterion revealed inconsistent results between supervisor and staff evaluations, where staff members rated higher than the supervisor, particularly for the complex prompt.

All in all, the findings suggest that \acp{llm} can generate suitable questions -- focusing content fidelity --, particularly when provided with a clear prompt and structured source material. For this task, DeepSeek and OpenAI were the best models, while Anthropic and Google were more variable in their performance. However, the low inter-rater reliability in the qualitative analysis significantly weakens the statistical power of the findings. Since the generated questions were not that challenging -- but more challenging with the complex prompt --, this reflects the need for Bloom's Taxonomy levels to be specified in the prompts, which can lead to higher \textit{Challenging} and \textit{Total Score} ratings.

\requ{2}{How does the relationship between diverse question formats and Bloom's Taxonomy levels influence the pedagogical effectiveness of \ac{llm}-generated questions?}

This research demonstrated that specifying Bloom levels -- either Bloom Only or Type+Bloom -- resulted in high adherence to the intended cognitive levels, with median scores of 10 in both cases. If no Bloom specification was given instead, the questions tended to fall into lower cognitive levels, with a median of 3 (at \textit{Understanding}), which is not a sufficient result for demanding educational content. The models performed diversely, with Anthropic and Google being rated higher than DeepSeek and OpenAI in Type Only conditions. However, when Bloom levels were specified, the performance of the models was higher, with OpenAI being the best model in the Bloom Only setting, and Google performing best in the Type+Bloom setting.

A critical finding was that Open-Ended questions consistently achieved higher cognitive levels than \acp{mcq}, and even higher Bloom alignment when a Bloom specification was given. The \acp{mcq} were more suitable for lower cognitive levels, such as \textit{Remembering}, \textit{Understanding}, and \textit{Applying}. By this, a mix of both question types, coupled with specifying Bloom levels -- which leads to a higher \textit{Total Score} -- can be generated to better handle all cognitive levels.

The student evaluation of the questions in Experiment 2 showed that the \textit{Bloom's Level} had the highest inter-rater reliability, supporting the usability of the clear level definitions used. However, linguistic and content-related criteria such as \textit{Relevance}, \textit{Clarity}, and \textit{Language} presented very low inter-rater reliability, as well as the \textit{Value} criterion. The evaluation rubrics used in this thesis were not sufficiently clear and consistent, which leads to the need for rubric refinements. Here, the current 0-10 scales provided too much interpretive freedom, and the students were not familiar with the ISO-OSI model domain. Overall, a balanced approach that combines Bloom's Taxonomy levels with appropriate question formats is essential for maximizing the pedagogical effectiveness of \ac{llm}-generated questions. Thoughtful prompt design, clear learning objectives, and a focus on linguistic clarity are crucial as well to ensure that the \ac{aqg} process is effective and beneficial for educational purposes, coupled with current \acp{llm}.

\subsection{Practical Implications for Education}

Firstly, valuable \ac{aqg} relies on concise, well-structured source materials (e.g.\ lecture scripts), which can produce more relevant and challenging questions than lengthy transcripts or book excerpts. Secondly, a well-designed prompting strategy is crucial for \ac{aqg}, as extensive prompt formulations do not necessarily lead to high content fidelity or overall quality. This was visible in Experiment 1a, where the common prompt outperformed the complex one in most quality dimensions, despite being less challenging. To achieve the best results, a concise but structured prompt including the necessary information, such as explicit Bloom levels or operator verbs, should be used. This helps to avoid overloading questions with unnecessary formulations or linguistic complexity, which can reduce \textit{Clarity} and \textit{Value}.

Overall, strategic selection of question formats for different cognitive levels is essential. Open-Ended questions are suitable across all Bloom levels and achieve superior cognitive alignment, while \acp{mcq} are constrained to lower levels, such as \textit{Remembering}, \textit{Understanding}, and \textit{Applying}. Therefore, a balanced mix of both question types can be generated to cover all cognitive levels, which is beneficial for exam preparation. A mixture of question type, Bloom level, source material, and learning objectives can be used to create a comprehensive question set.

Moreover, users, particularly students, should be aware that \acp{llm} are not perfect and can produce hallucinations or incorrect information. Therefore, it is important to inform students about the limitations of \acp{llm}, and integrate a feedback mechanism into future frameworks, where students can provide feedback on the quality of the generated questions. This can help improve the quality of the questions over time, and ensure that they meet the needs of the students.

\subsection{Future Research Directions}

The dynamic landscape of \acp{llm} presents both challenges and opportunities for future research. As newer models are going to be released, they may exhibit even greater potential for \ac{aqg} and educational applications, while mitigating existing limitations. Hallucinations remain a critical issue, since factual accuracy is a crucial concern. Future work should continue to reduce the generation of incorrect or misleading information, especially in educational contexts.

\pagebreak

Based on this, analyzing the adherence of \ac{llm}-generated questions to educational content can be further improved by using a more comprehensive multi-\ac{llm} judging approach, where the questions are better analyzed by multiple models instead of just two and averaging both results, which was done in this thesis. This would imply a more robust foundation for automated evaluation without solely relying on expert judgments, as these are often time-consuming. To avoid randomness, which might occur in the evaluation process above, more recent or future embedding models should be utilized instead of a sentence transformer, which was used in this thesis. Since sentence transformers are not the best choice for semantic similarity in longer texts\,/\,varying structures, this could be improved by addressing this issue in future work.

The experimental setup for analyzing \textit{Manipulation Handling} was problematic, with a sample size of only 16/56 questions being evaluated qualitatively. A larger evaluation set would provide more consistent insights into the \acp{llm}' reaction to manipulations, as the results were inconsistent between the supervisor and staff members. Additionally, the quantitative analysis using all 56 questions showed different trends than the qualitative analysis, highlighting the sampling limitations for this sub-experiment. This also implies a redesign of the rubric, since the \textit{Value} criterion was criticized for its lack of explicit learning objectives integration, which was not considered during the planning phase. Furthermore, clearer anchor descriptions -- e.g.\ for 0 points, 2 points, ... -- are needed to improve inter-rater reliability and thus the qualitative analysis of the questions.

While this thesis focused on \textit{reasoning} models, future research should also compare \textit{reasoning} and non-\textit{reasoning} models to analyze differences in \ac{aqg}. Unfortunately, the \textit{reasoning} step before answering the user is often not well analyzable, as many models consider this step proprietary. Based on this, the developers are solely providing summaries of the \textit{reasoning} process instead\footnote{\url{https://platform.openai.com/docs/guides/reasoning?api-mode=responses\#reasoning-summaries}, accessed 08/12/2025}.

Generating \acp{mcq} in one step is still a challenge, as in few cases, the (in)correct answer options are not generated properly, which can lead to confusion for students. Based on this, a pipeline that generates the question, correct answer(s), and incorrect answer options separately would be a beneficial approach for future work, as it was not implemented in this thesis, but demonstrated in work such as \cite{bhowmick_automating_2023}. Finally, addressing potential biases in the evaluation process is crucial, as the supervisor of this thesis was also the author of the script/transcript used for question generation. However, the overall trends in the evaluation were similar to those of the other raters. In addition, the number of raters was sufficient for this thesis (five raters in Experiment 1 and three raters in Experiment 2), but the poor inter-rater reliability -- especially in Experiment 1 -- indicates that increasing the number of raters alone would not resolve fundamental issues with rubric design. This aspect should be addressed in future work, to enhance the robustness of the findings. All these insights can be used to build a robust, reliable and particularly pedagogical framework for \ac{aqg}, which is a promising direction for further research.